\chapter{State-of-the-Art and Progress Beyond} \label{ch:state_of_the_art}

This chapter summarizes the current state of knowledge in the literature regarding BESS allocation and advanced energy system modeling. Section \ref{sec:soa_bess} provides a comprehensive review of BESS allocation in power grids, utilizing a multi-dimensional taxonomy to identify the prevailing research gap between strictly technical grid simulations and purely economic optimization approaches. Section \ref{sec:soa_epec} explores the application of EPECs in modern energy economics, highlighting their capability to realistically capture strategic, oligopolistic market behavior. Finally, Section \ref{sec:novelty} synthesizes these findings to clearly define the genuine contribution of this thesis and outline the progress beyond the current state-of-the-art.

\section{State-of-the-Art: BESS Allocation in Power Grids} \label{sec:soa_bess}
The allocation of BESS in power grids has been extensively studied from various methodological perspectives. By classifying recent literature through the multi-dimensional taxonomy in Table \ref{tab:bess_summary_taxonomy} -- focusing on optimization methods, objective functions, and constraints -- distinct trends and research gaps emerge. 

\begin{table}[ht]
\centering
\caption{Taxonomy of BESS allocation in power grids. (Aug.) refers to battery augmentation, whereas (Dist.) refers to distribution across existing generators.}
\label{tab:bess_summary_taxonomy}
\resizebox{\linewidth}{!}{%
\begin{tabular}{|l|c|c|l|p{3.2cm}|p{3.2cm}|p{3.0cm}|p{2.0cm}|p{2.2cm}|}
\hline
\textbf{Literature} & \textbf{Sizing} & \textbf{Siting} & \textbf{Focus} & \textbf{Optimization Method} & \textbf{Objective function}& \textbf{Constraints}& \textbf{Time Horizon} & \textbf{Local Focus} \\
\hline
\hline
\cite{al_ahmad_optimal_2020} & X & X & Both & Hybrid MOPSO \& NSGAII & Min. Cost, Voltage Deviation, CO\textsubscript{2} & Grid Simulation & Today & Generic (IEEE 30-bus) \\
\hline
\cite{das_optimal_2020} & X & & Technical & FSCABC, verified with PSO & Min. Freq. Deviation \& ROCOF & Grid Sim (Freq. Nadir, ROCOF) & Today/Near-Future & Generic (IEEE 39-bus) \\
\hline
\cite{shin_optimal_2020} & X (Aug.) & & Economical & Iterative Brute-Force Search & Max. Net Present Value (NPV) & Battery Degradation, Market Sim (PPA) & Future (15-year) & South Korea \\
\hline
\cite{wang_optimal_2020} & X & X & Technical & Improved PSO (IPSO) & Min. Voltage Deviation \& Power Loss (Fuzzy) & Grid Simulation & Today & Generic (IEEE-33 bus) \\
\hline
\cite{hameed_business-oriented_2021} & & X & Economical & Qualitative Scoring Framework & Max. Business Feasibility & External/Business Factors (Land, Permits) & Today/Near-Future & Denmark (Bornholm) \\
\hline
\cite{mohamad_optimum_2021} & X & X & Technical & GA with Monte Carlo Sim & Min. BESS Units \& Solar Curtailment & Grid Simulation (DC-OPF) & Today & Generic (IEEE 24-bus) \\
\hline
\cite{odero_wind_2022} & X & & Economical & BPNN (prediction) + GA (opt.) & Min. Investment Cost & Reliability (LPSP), Grid Sim & Today/Near-Future & Kenya (LTWPP) \\
\hline
\cite{wang_energy_2022} & X & X & Technical & Lssvm (partition) + ResNet DNN & Max. Energy Throughput & Grid Simulation & Today & Generic (IEEE 33-node) \\
\hline
\cite{alonso_sorensen_energy_2023} & X (Dist.) & & Technical & Analytical (Swing Equation) & Limit ROCOF per Area & Local Frequency Dynamics & Future & Generic \& Cape Verde \\
\hline
\cite{czock_place_2023} & & X & Economical & Linear Optimization (LOPF) & Min. Total System Costs & Grid Simulation (DC-OPF) & Future (2030) & Germany \\
\hline
\end{tabular}
}
\end{table}

Regarding the \textbf{optimization methods}, the literature is broadly divided into meta-heuristic approaches and deterministic models. Meta-heuristic algorithms, such as Genetic Algorithms (GA) or Particle Swarm Optimization (PSO), are frequently used to navigate the non-linearities of BESS placement and sizing \cite{das_optimal_2020}. Conversely, deterministic approaches utilize Linear Optimal Power Flow (LOPF) to find guaranteed optimal solutions for specific system setups, often minimizing total system costs \cite{czock_place_2023}. Other works employ highly intensive brute-force searches to evaluate complex, non-linear battery degradation models for BESS sizing \cite{shin_optimal_2020}.

When analyzing the \textbf{objective functions and constraints}, a significant dichotomy becomes apparent. Most studies focus strictly on either technical or economic goals. Technical models prioritize grid stability, utilizing BESS to minimize frequency deviations, Rate of Change of Frequency (ROCOF), or voltage fluctuations \cite{das_optimal_2020, alonso_sorensen_energy_2023}. In these models, economic realities like investment costs are often heavily simplified. On the other hand, purely economic models aim to maximize the Net Present Value (NPV) through energy arbitrage, but often neglect strict physical grid constraints (such as nodal congestion) or the physical degradation of the battery assets \cite{shin_optimal_2020}. 

Ultimately, the current state-of-the-art lacks a comprehensive, hybrid approach. There is a distinct gap for optimization frameworks that seamlessly combine detailed economic profit maximization (such as multi-service market participation in both spot and reserve markets) with strict technical and physical grid constraints.

\section{State-of-the-Art: EPECs in Energy Systems Modeling} \label{sec:soa_epec}
While traditional cost-minimization models (like central-planner LOPF) are valuable for theoretical grid studies, they fail to represent the reality of liberalized electricity markets, where independent actors maximize their own profits. As emphasized by Conejo and Ruiz, ``Complementarity, not Optimization, is the Language of Markets'' \cite{conejo_complementarity_2020}. To accurately model these strategic interactions, the methodology of Equilibrium Problems with Equilibrium Constraints (EPECs) has gained significant traction.

An EPEC finds the Nash Equilibrium among multiple strategic actors, where each actor's problem is formulated as a Mathematical Program with Equilibrium Constraints (MPEC) \cite{conejo_complementarity_2020}. This methodology has proven highly effective in modern energy economics. For instance, Chen et al. successfully applied an EPEC to determine equilibria in coupled electricity and natural gas markets, demonstrating the method's ability to capture complex strategic offering and bidding behaviors \cite{chen_equilibria_2020}. 

Particularly relevant to this thesis is the work by Wogrin et al., who utilized an EPEC to model generation capacity expansion in liberalized markets \cite{wogrin_capacity_2013}. They proved that a bilevel approach -- where investors maximize profits in the upper level while anticipating the market clearing in the lower level -- yields vastly different and more realistic investment schedules compared to traditional single-stage cost-minimization models. However, while EPECs are established for conventional generation expansion, their application to the highly complex, constraint-heavy allocation of multi-use BESS remains a largely unexplored frontier.

\section{Contribution} \label{sec:novelty}
Based on the identified gaps in BESS allocation and the proven strengths of complementarity modeling, this thesis creates a novel methodological bridge. The genuine contribution of this work lies in the development of a comprehensive EPEC framework tailored specifically to the multi-service allocation of BESS. 

The progress beyond the state-of-the-art is defined by the following key innovations:
\begin{itemize}
    \item \textbf{Decentralized Oligopoly Modeling:} Shifting away from central-planner assumptions by modeling strategic BESS investors who actively compete for nodal grid capacities and market shares.
    \item \textbf{Simultaneous Siting and Sizing:} Endogenously determining both the optimal location and the exact Energy-to-Power ratio (MW and MWh) based on individual investor risk affinities (WACC).
    \item \textbf{Multi-Service Co-Optimization:} Bridging the gap between technical and economic models by co-optimizing day-ahead spot market arbitrage with the strict physical requirements of continuous reserve market provision (aFRR).
    \item \textbf{Internalized Physical Constraints:} Integrating linearized cyclic battery degradation directly into the strategic objective function, ensuring that investment decisions reflect long-term operational wear.
\end{itemize}

Based on the identified gaps in BESS allocation and the proven strengths of complementarity modeling, this thesis creates a novel methodological bridge. The genuine contribution of this work lies in the development of a comprehensive EPEC framework tailored specifically to the multi-service allocation of BESS. So far, the literature has predominantly focused on either purely technical grid simulations or centralized economic cost-minimization models, while the application of strategic equilibrium models has largely been restricted to conventional generation expansion.

In contrast, this thesis shifts away from traditional central-planner assumptions by introducing a decentralized oligopoly model representing strategic BESS investors who actively compete for nodal grid capacities and market shares. Rather than relying on simplified heuristics, the framework simultaneously determines the optimal location and the exact Energy-to-Power ratio (MW and MWh) based on individual investor risk affinities. Furthermore, it bridges the gap between technical and economic models through a multi-service co-optimization that balances day-ahead spot market arbitrage with the strict physical requirements of continuous reserve market provision (aFRR). Ultimately, this approach fully internalizes physical constraints by integrating linearized cyclic battery degradation directly into the strategic objective function, ensuring that the resulting investment decisions realistically reflect long-term operational wear.


